%% kapa: exam=1 round=21 year=2010 period=1 form=A subject=economics subject_ko=경제원론 seq=27 number=67
%% source: Q-Net 감정평가사 기출문제 (cst00309) artlSeq=5024181 file=extracted/1008/5024181/2010년 제21회 감정평가사 1차시험 기출문제/1교시.hwp
%% conversion: kordoc 4.24.51 -> tools/gen1.py (auto); review=proofread-sonnet; flags=none
\begin{question}{67}
달러화에 대한 원화의 실질환율과 명목환율에 관한 설명으로 옳지 않은 것은?
\begin{choices}
  \item 명목환율이 일정할 때 실질환율이 상승(절하)되면 미국 제품에 비해 우리나라 제품의 가격이 더 비싸진다.
  \item 양국의 물가수준이 일정할 때 명목환율이 상승(절하)하면 실질환율도 상승(절하)한다.
  \item 실질환율이 하락(절상)되면 장기적으로 우리나라의 순수출은 감소한다.
  \item 구매력평가설에 따르면 미국의 물가수준이 상승하고 우리나라 물가수준이 하락할 때 명목환율이 변한다.
  \item 구매력평가설에 따르면 우리나라 제품의 가격과 미국 제품의 원화표시 가격의 상대적 비율은 일정하다.
\end{choices}
\end{question}
\answer{1}
